\documentclass[11pt]{article}

\usepackage[margin=1in]{geometry}
\usepackage{booktabs}
\usepackage{multirow}
\usepackage{float}
\usepackage{amsmath}
\usepackage{microtype}
\usepackage{authblk}
\usepackage[hidelinks]{hyperref}
\usepackage{xspace}

\newcommand{\B}[1]{\textbf{#1}}

\newcommand{\gfaz}{GFAz\xspace}

\setlength{\parskip}{4pt}
\setlength{\parindent}{0pt}

\title{\vspace{-2.5em}\gfaz: Order-of-Magnitude Pangenome Analytics by Computing over Compression}

\author[1]{Taolue Yang}
\author[1]{Youyuan Liu}
\author[1]{Bo Jiang}
\author[1]{Xinghua Shi}
\author[1]{Sian Jin}
\affil[1]{Temple University, Philadelphia, PA, USA}
\date{}

\begin{document}
\maketitle
\vspace{-3em}

\paragraph{Overview.}
Pangenome graphs in the Graphical Fragment Assembly (GFA) format now reach hundreds
of gigabytes, with haplotype traversal records exceeding 90\% of file size (98.7\% of
a 14\,GB graph of 1{,}000 chr19 haplotypes). This imposes two costs. For storage,
general-purpose compressors (gzip, zstd) work over local byte windows and miss the
cross-haplotype redundancy in traversals, while specialized tools drop GFA semantics
or offer low throughput. For analysis, GFA is the format pangenomes are built in but
rarely analyzed in: downstream tools (bcftools, GATK, plink, GWAS) consume
reference-anchored tables such as VCF, so each task first materializes a whole-graph
object or snarl index larger than the statistic it computes. \gfaz addresses both
with one artifact: a lossless GFA compressor whose container doubles as a compute
substrate, running path-centric analyses directly on the grammar-compressed
traversal store.

\paragraph{Design.}
\emph{Compression.} \gfaz normalizes paths and walks to signed \texttt{int32}
traversals (sign = strand) and applies multi-round 2-mer grammar compression that
exploits forward/reverse orientation symmetry: a negative rule reference denotes a
rule's reverse complement, so one rulebook serves both orientations. All record
types, including optional fields, are stored as type-specific columns, giving fully
lossless round-trips, per-field access, incremental haplotype extension, and
single-haplotype random access. Both compression and decompression are linear-time
and fully parallel, with interchangeable OpenMP (CPU) and CUDA/nvComp (GPU) backends
that share one format.
\emph{Compute engine.} A streaming primitive expands one traversal to
orientation-signed node IDs straight from the grammar (array-indexed rule lookup, no
hashing) and folds it into a compact accumulator, so a query touches only the columns
it needs and peak memory stays below the uncompressed GFA. The flagship analysis
\texttt{deconstruct} (GFA to VCF) is a new two-phase algorithm that separates
one-time, topology-only site discovery (biconnected decomposition of the node-end
graph, $O(V{+}E)$) from a streaming allele-observation fold, and never builds a graph
or a snarl index. Five further analyses (\texttt{growth}, \texttt{pav},
\texttt{depth}, \texttt{similarity}, \texttt{stats}) reuse the same substrate, each
reproducing an established tool's output.

\paragraph{Compression performance.}
All results use a single node with an AMD Ryzen Threadripper PRO 9955WX (16 cores,
32 threads), an NVIDIA RTX Pro 6000 GPU, and 512\,GB of DDR5-6400 memory (503\,GiB
usable). Across six datasets (113\,MB to 369\,GB), \gfaz gives the best ratio of
every method on both backends, 18 to 84$\times$, typically over 10$\times$ better
than gzip/zstd and ahead of the specialized GBZ and sqz; the best sqz variant trails
\gfaz by about 2$\times$ at chromosome scale and exhausts memory on every HPRC graph.
It also leads compression and decompression throughput where inputs fit: the CPU
backend sustains up to 385\,MB/s and 1.8\,GB/s, and the GPU backend up to 4.8\,GB/s
and 9.4\,GB/s. The ratio advantage widens with graph complexity, reaching
83.9$\times$ on HPRC v2.0 versus 66.8$\times$ (GBZ) and 6.5$\times$ (zstd)
(Table~\ref{tab:compress}).

\begin{table}[H]
\centering
\scriptsize
\caption{Compression ratio and throughput (MB/s) vs.\ baselines. Ratio = ×; Co. =
compression speed, De. = decompression speed. \B{Bold} = best in row. ``sqz'' is the
best-ratio sqz variant (sqz+bgzip); it runs out of memory on all HPRC graphs (NA).
GPU throughput on the two largest graphs is omitted (traversals exceed device memory;
ratio shown via CPU simulation). GBZ and sqz preserve fewer GFA semantics.}
\label{tab:compress}
\setlength{\tabcolsep}{3pt}
\begin{tabular}{llrrrrrr}
\toprule
Dataset (GFA) & & Gzip & Zstd & sqz & GBZ & \gfaz(CPU) & \gfaz(GPU) \\
\midrule
\multirow{3}{*}{chr1 (7.1\,GB)}    & Ratio & 5.59 & 7.54 & 18.0 & 9.52 & \B{35.4} & 31.7 \\
                                   & Co.   & 46   & 2178 & 4.0  & 12   & 385  & \B{2754} \\
                                   & De.   & 359  & 1618 & 21   & 284  & 1658 & \B{8124} \\
\midrule
\multirow{3}{*}{chr6 (3.8\,GB)}    & Ratio & 5.04 & 6.99 & 20.8 & 19.2 & \B{35.4} & 28.2 \\
                                   & Co.   & 41   & 1712 & 3.6  & 11   & 348  & \B{3791} \\
                                   & De.   & 348  & 1515 & 20   & 281  & 1758 & \B{7230} \\
\midrule
\multirow{3}{*}{\emph{E. coli} (113\,MB)} & Ratio & 4.69 & 5.67 & 7.46 & 5.58 & \B{18.3} & 16.7 \\
                                   & Co.   & 33   & \B{1356} & 4.5 & 20 & 190 & 678 \\
                                   & De.   & 310  & 1258 & 32  & 197 & 491  & \B{1430} \\
\midrule
\multirow{3}{*}{HPRC v1.1 (48\,GB)} & Ratio & 4.02 & 5.32 & NA & 14.0 & \B{22.4} & 20.4 \\
                                   & Co.   & 36   & 1657 & NA & 85   & 231  & \B{4843} \\
                                   & De.   & 319  & 1234 & NA & 650  & 1058 & \B{9435} \\
\midrule
\multirow{3}{*}{HPRC v2.0 (358\,GB)} & Ratio & 4.19 & 6.49 & NA & 66.8 & \B{83.9} & 76.4 \\
                                   & Co.   & 49   & \B{1514} & NA & 130 & 367 & n/a \\
                                   & De.   & 342  & 1240 & NA & 648  & \B{1652} & n/a \\
\midrule
\multirow{3}{*}{HPRC v2.1 (369\,GB)} & Ratio & 4.19 & 6.43 & NA & 64.2 & \B{82.8} & 74.2 \\
                                   & Co.   & 49   & \B{1540} & NA & 136 & 348 & n/a \\
                                   & De.   & 343  & 1241 & NA & 652  & \B{1559} & n/a \\
\bottomrule
\end{tabular}
\end{table}

\paragraph{Compute performance.}
Computing on the container makes \gfaz one to three orders of magnitude faster than
the standard tools at a fraction of their memory, while matching their output:
\texttt{deconstruct} agrees with \texttt{vg deconstruct} at about 99.99\% of
reference positions, \texttt{growth} reproduces Panacus curves exactly, and
\texttt{pav} reproduces \texttt{odgi pav} semantics (Table~\ref{tab:downstream}). The
sub-GFA memory footprint also lets \gfaz scale to whole human pangenomes the
baselines cannot load: from a 4.5\,GB container it emits the genome-wide VCF (about
35\,M records) on the 358 and 369\,GB HPRC v2 graphs in about 40\,min at 191 to
195\,GB peak RSS (roughly half the GFA), and the exact growth curve in about 40\,s at
only 13 to 14\,GB (about 27$\times$ smaller than the GFA).

\begin{table}[H]
\centering
\scriptsize
\caption{Downstream analytics vs.\ standard tools (16 threads): wall-clock time and
peak RSS. \gfaz reads the compressed \texttt{.gfaz}; baselines read the GFA. ``DNF''
= baseline does not complete in practical memory; on whole-genome HPRC graphs \gfaz
runs alone. (odgi \texttt{pav} additionally needs an \texttt{odgi build} step, 31--61\,s;
whole-genome \texttt{pav} on the two largest graphs exceeds the 503\,GB node and is
omitted.)}
\label{tab:downstream}
\setlength{\tabcolsep}{4pt}
\begin{tabular}{llrrrrr}
\toprule
 & & \multicolumn{2}{c}{\gfaz} & \multicolumn{2}{c}{Baseline} & \\
\cmidrule(lr){3-4}\cmidrule(lr){5-6}
Analysis & Graph (GFA) & Time & RSS & Time & RSS & Speedup \\
\midrule
\multirow{5}{*}{\shortstack[l]{\texttt{deconstruct}\\\footnotesize vs.\ \texttt{vg}}}
  & chr1 (7.1\,GB)   & 11.3\,s   & 8.3\,GB  & 805\,s & 30.3\,GB & 71$\times$ \\
  & chr6 (3.8\,GB)   & 7.0\,s    & 5.2\,GB  & 349\,s & 18.4\,GB & 50$\times$ \\
  & HPRC v1.1 (48\,GB)  & 5.1\,min  & 37.4\,GB & DNF & n/a & n/a \\
  & HPRC v2.0 (358\,GB) & 39.8\,min & 191\,GB  & DNF & n/a & n/a \\
  & HPRC v2.1 (369\,GB) & 41.4\,min & 195\,GB  & DNF & n/a & n/a \\
\midrule
\multirow{5}{*}{\shortstack[l]{\texttt{growth}\\\footnotesize vs.\ Panacus}}
  & chr1 (7.1\,GB)   & 0.74\,s & 0.66\,GB & 18.3\,s & 6.16\,GB & 25$\times$ \\
  & chr6 (3.8\,GB)   & 0.36\,s & 0.31\,GB & 9.0\,s  & 3.77\,GB & 25$\times$ \\
  & HPRC v1.1 (48\,GB)  & 6.2\,s  & 6.3\,GB  & DNF & n/a & n/a \\
  & HPRC v2.0 (358\,GB) & 39.8\,s & 12.9\,GB & DNF & n/a & n/a \\
  & HPRC v2.1 (369\,GB) & 41.2\,s & 13.7\,GB & DNF & n/a & n/a \\
\midrule
\multirow{3}{*}{\shortstack[l]{\texttt{pav}\\\footnotesize vs.\ \texttt{odgi}}}
  & chr1 (7.1\,GB)   & 11.7\,s  & 9.5\,GB & 3499\,s & 31.9\,GB & 299$\times$ \\
  & chr6 (3.8\,GB)   & 7.8\,s   & 6.4\,GB & 4759\,s & 19.4\,GB & 613$\times$ \\
  & HPRC v1.1 (48\,GB)  & 1.8\,min & 77\,GB  & DNF & n/a & n/a \\
\bottomrule
\end{tabular}
\end{table}

\paragraph{Conclusion.}
\gfaz folds the two recurring costs of large pangenomes, storage and derivation, into
a single container that is state-of-the-art at compression and a fast, low-memory
engine for the analyses practitioners run in practice. It is open source at
\url{https://github.com/babyplutokurt/GFAz}.

\end{document}
